\sgaStableTarget{sga6.e07.prop.3-4}\textbf{Proposition 3.4.} \uline{With the hypotheses and notation of} 3.1 \uline{and} 3.3, \uline{for every} @@M0@@ \uline{we have the relation}:
\sgaStableTarget{sga6.e07.eq.3-4-1}@@M1@@

Since both sides of the relation to be proved are additive, we may suppose that @@M2@@ is the class of an @@M3@@-Module locally free of finite type @@M4@@. The element @@M5@@ is then the class of the complex @@M6@@, obtained by taking a projective resolution of @@M7@@ and applying the functor @@M8@@ to it. The homology sheaves of the resulting complex are the @@M9@@. By 2.4 and 3.2, we have:
@@M10@@

It follows that these homology sheaves are locally free of finite type over @@M11@@, hence perfect over @@M12@@; consequently, the class in @@M13@@ of the complex @@M14@@ is the alternating sum of the classes of its homology sheaves, namely:
@@M15@@

\sgaStableTarget{sga6.e07.lem.3-5}\textbf{Lemma 3.5.} \uline{With the hypotheses and notation of} 3.1, \uline{we have the following isomorphisms}:
\begin{enumerate}[label=\roman*)]
\item \uline{for every} @@M16@@, \uline{and every} @@M17@@, @@M18@@;
\item \uline{for every} @@M19@@, @@M20@@.
\end{enumerate}

We may clearly suppose @@M21@@ affine, hence quasi-compact.

i) First recall the following lemma, which was proved in exp. III.
